% i18n_qa_hi_paper.tex — Hindi statistics paper on Bayesian inference
% Intentional errors: ZWJ misuse between conjuncts, mixing Hindi and
%   Arabic numerals, Devanagari quotation marks issues
\documentclass[12pt,a4paper]{article}

\usepackage[hindi]{babel}
\usepackage[utf8]{inputenc}
\usepackage[T1]{fontenc}
\usepackage{amsmath}
\usepackage{amssymb}
\usepackage{amsthm}
\usepackage{graphicx}
\usepackage{hyperref}
\usepackage[margin=2.5cm]{geometry}
\usepackage{booktabs}
\usepackage{natbib}
\usepackage{enumitem}
\usepackage{microtype}
\usepackage{xcolor}
\usepackage{float}
\usepackage{mathrsfs}
\usepackage{bm}
\usepackage{tikz}
\usepackage{fontspec}

\theoremstyle{plain}
\newtheorem{theorem}{प्रमेय}[section]
\newtheorem{lemma}[theorem]{लेम्मा}
\newtheorem{proposition}[theorem]{प्रस्ताव}
\newtheorem{corollary}[theorem]{उपप्रमेय}

\theoremstyle{definition}
\newtheorem{definition}[theorem]{परिभाषा}
\newtheorem{example}[theorem]{उदाहरण}

\theoremstyle{remark}
\newtheorem{remark}[theorem]{टिप्पणी}

\title{बायेसियन अनुमान के लिए मार्कोव श्रृंखला\\
मोंटे कार्लो विधियों का तुलनात्मक अध्ययन}
\author{%
  प्रो॰ अमित शर्मा\thanks{सांख्यिकी विभाग, भारतीय सांख्यिकी संस्थान, कोलकाता।}
  \and
  डॉ॰ प्रिया गुप्ता\thanks{गणित विभाग, भारतीय प्रौद्योगिकी संस्थान, दिल्ली।}
}
\date{६ अप्रैल २०२६}

\begin{document}

\maketitle

\begin{abstract}
इस शोधपत्र में हम बायेसियन अनुमान के लिए विभिन्न मार्कोव श्रृंखला
मोंटे कार्लो (MCMC) विधियों का तुलनात्मक अध्ययन प्रस्तुत करते हैं।
हम गिब्स सैम्पलर, मेट्रोपोलिस-हेस्टिंग्स एल्गोरिदम, और हैमिल्टोनियन
मोंटे कार्लो (HMC) विधि की तुलना उच्च-आयामी पश्च वितरणों पर करते हैं।
% INTENTIONAL ERROR: mixing Hindi numerals (१२) and Arabic numerals (50)
हमारे प्रयोगात्मक परिणाम दर्शाते हैं कि HMC विधि १२ से 50 आयामों वाले
प्रतिदर्श स्थानों में अभिसरण दर में ३-५ गुना सुधार प्रदान करती है।
हम एक नया अनुकूली HMC एल्गोरिदम भी प्रस्तुत करते हैं जो स्वचालित रूप
से पथ लंबाई और चरण आकार को समायोजित करता है। सैद्धांतिक विश्लेषण
से पता चलता है कि प्रस्तावित विधि सभी लॉग-अवतल लक्ष्य वितरणों के लिए
ज्यामितीय एर्गोडिकता प्राप्त करती है।
\end{abstract}

\tableofcontents

%% ---------------------------------------------------------------
\section{परिचय}
\label{sec:intro}

बायेसियन सांख्यिकी में पश्च वितरण का अनुमान एक केंद्रीय समस्या है।
बायेस प्रमेय के अनुसार, पश्च वितरण इस प्रकार दिया जाता है:
\begin{equation}
\label{eq:bayes}
p(\bm{\theta} \mid \mathbf{y}) = \frac{p(\mathbf{y} \mid \bm{\theta}) \, p(\bm{\theta})}{p(\mathbf{y})},
\end{equation}
जहाँ $\bm{\theta}$ पैरामीटर सदिश है, $\mathbf{y}$ प्रेक्षित आँकड़े हैं,
$p(\mathbf{y} \mid \bm{\theta})$ प्रशंसनीयता फलन है, $p(\bm{\theta})$ पूर्व
वितरण है, और $p(\mathbf{y}) = \int p(\mathbf{y} \mid \bm{\theta}) p(\bm{\theta}) \, d\bm{\theta}$
सीमांत प्रशंसनीयता है।

% INTENTIONAL ERROR: ZWJ (‍) misused between conjuncts — unnecessary ZWJ in सम्‍भव
अधिकांश व्यावहारिक मामलों में, सीमांत प्रशंसनीयता का विश्लेषणात्मक
मूल्यांकन सम्‍भव नहीं होता\footnote{केवल संयुग्मी पूर्व वितरण परिवारों
के लिए बंद-रूप समाधान उपलब्ध हैं।}। इसलिए, MCMC विधियाँ पश्च
वितरण से प्रतिदर्श उत्पन्न करने का मानक उपकरण बन गई हैं~\citep{gelman2013}.

\subsection{मार्कोव श्रृंखला मोंटे कार्लो}

MCMC विधियाँ एक मार्कोव श्रृंखला $\{\bm{\theta}^{(t)}\}_{t=0}^{T}$
का निर्माण करती हैं जिसका स्थिर वितरण लक्ष्य पश्च वितरण
$p(\bm{\theta} \mid \mathbf{y})$ है।

\begin{definition}
\label{def:ergodic}
एक मार्कोव श्रृंखला ज्यामितीय रूप से एर्गोडिक है यदि कोई
$\rho \in (0,1)$ और $C > 0$ अस्तित्व में हैं जैसे कि
\begin{equation}
\label{eq:ergodic}
\|P^t(\bm{\theta}, \cdot) - \pi(\cdot)\|_{\mathrm{TV}} \leq C \, \rho^t,
\end{equation}
जहाँ $P^t$ संक्रमण नाभिक की $t$-वीं शक्ति है और $\pi$ स्थिर
वितरण है।
\end{definition}

% INTENTIONAL ERROR: mixing Arabic numerals (1953) with Hindi text
मेट्रोपोलिस एल्गोरिदम सर्वप्रथम 1953 में प्रस्तावित किया गया
था~\citep{metropolis1953} और बाद में हेस्टिंग्स (1970) द्वारा
सामान्यीकृत किया गया~\citep{hastings1970}।

\subsection{शोधपत्र की संरचना}

% INTENTIONAL ERROR: Devanagari quotation marks — using English "..." instead of proper marks
इस शोधपत्र की संरचना निम्न प्रकार है: खंड~\ref{sec:methods} में
हम "मानक MCMC विधियों" का वर्णन करते हैं, खंड~\ref{sec:hmc}
में हम HMC विधि और हमारे अनुकूली एल्गोरिदम का विवरण देते हैं,
खंड~\ref{sec:theory} में सैद्धांतिक गुणधर्म प्रस्तुत करते हैं,
और खंड~\ref{sec:experiments} में प्रयोगात्मक परिणाम दिखाते हैं।

%% ---------------------------------------------------------------
\section{मानक MCMC विधियाँ}
\label{sec:methods}

\subsection{मेट्रोपोलिस-हेस्टिंग्स एल्गोरिदम}

मेट्रोपोलिस-हेस्टिंग्स (MH) एल्गोरिदम प्रत्येक चरण में एक
प्रस्ताव $\bm{\theta}^*$ उत्पन्न करता है और इसे स्वीकृति
प्रायिकता के साथ स्वीकार या अस्वीकार करता है:
\begin{equation}
\label{eq:mh_accept}
\alpha(\bm{\theta}^{(t)}, \bm{\theta}^*) = \min\!\left(1, \frac{p(\bm{\theta}^* \mid \mathbf{y}) \, q(\bm{\theta}^{(t)} \mid \bm{\theta}^*)}{p(\bm{\theta}^{(t)} \mid \mathbf{y}) \, q(\bm{\theta}^* \mid \bm{\theta}^{(t)})}\right),
\end{equation}
जहाँ $q(\cdot \mid \cdot)$ प्रस्ताव वितरण है।

\begin{proposition}
\label{prop:mh_balance}
MH एल्गोरिदम विस्तृत संतुलन शर्त को संतुष्ट करता है:
\[
\pi(\bm{\theta}) \, P(\bm{\theta}, \bm{\theta}') = \pi(\bm{\theta}') \, P(\bm{\theta}', \bm{\theta}),
\]
जहाँ $\pi = p(\cdot \mid \mathbf{y})$ लक्ष्य वितरण है।
\end{proposition}

\subsection{गिब्स सैम्पलर}

गिब्स सैम्पलर MH का एक विशेष मामला है जहाँ प्रस्ताव
सशर्त वितरणों से लिया जाता है। $d$-आयामी पैरामीटर
$\bm{\theta} = (\theta_1, \ldots, \theta_d)$ के लिए:

% INTENTIONAL ERROR: mixing Hindi numerals ३ and Arabic numerals 1
चरण ३.1: प्रत्येक $i = 1, \ldots, d$ के लिए:
\begin{equation}
\label{eq:gibbs}
\theta_i^{(t+1)} \sim p(\theta_i \mid \theta_1^{(t+1)}, \ldots, \theta_{i-1}^{(t+1)}, \theta_{i+1}^{(t)}, \ldots, \theta_d^{(t)}, \mathbf{y}).
\end{equation}

\begin{remark}
गिब्स सैम्पलर को सभी सशर्त वितरणों का ज्ञान आवश्यक है,
जो व्यवहार में हमेशा उपलब्ध नहीं होता\footnote{%
संयुग्मी मॉडलों में सशर्त वितरण ज्ञात रूपों में आते हैं,
जैसे सामान्य-सामान्य या बीटा-द्विपद।}।
\end{remark}

%% ---------------------------------------------------------------
\section{हैमिल्टोनियन मोंटे कार्लो}
\label{sec:hmc}

\subsection{भौतिकी प्रेरणा}

HMC विधि हैमिल्टोनियन यांत्रिकी से प्रेरित है। सहायक
संवेग चर $\mathbf{p}$ प्रवर्तित करते हुए, हैमिल्टोनियन
फलन परिभाषित करें:
\begin{equation}
\label{eq:hamiltonian}
\mathcal{H}(\bm{\theta}, \mathbf{p}) = -\log p(\bm{\theta} \mid \mathbf{y}) + \frac{1}{2} \mathbf{p}^\top M^{-1} \mathbf{p},
\end{equation}
जहाँ $M$ द्रव्यमान आव्यूह है। हैमिल्टन समीकरण हैं:
\begin{align}
\label{eq:hamilton_eqs}
\frac{d\bm{\theta}}{dt} &= M^{-1} \mathbf{p}, \\
\frac{d\mathbf{p}}{dt} &= \nabla_{\bm{\theta}} \log p(\bm{\theta} \mid \mathbf{y}).
\end{align}

\subsection{लीपफ्रॉग समाकलक}

हैमिल्टन समीकरणों का संख्यात्मक समाकलन लीपफ्रॉग विधि
से किया जाता है। चरण आकार $\varepsilon$ और $L$ चरणों के लिए:
\begin{align}
\mathbf{p}_{1/2} &= \mathbf{p}_0 + \frac{\varepsilon}{2} \nabla_{\bm{\theta}} \log p(\bm{\theta}_0 \mid \mathbf{y}), \label{eq:leapfrog1} \\
\bm{\theta}_1 &= \bm{\theta}_0 + \varepsilon \, M^{-1} \mathbf{p}_{1/2}, \label{eq:leapfrog2} \\
\mathbf{p}_1 &= \mathbf{p}_{1/2} + \frac{\varepsilon}{2} \nabla_{\bm{\theta}} \log p(\bm{\theta}_1 \mid \mathbf{y}). \label{eq:leapfrog3}
\end{align}

\subsection{अनुकूली HMC एल्गोरिदम}

हम प्रस्ताव करते हैं कि चरण आकार $\varepsilon$ और पथ लंबाई $L$
को स्वचालित रूप से समायोजित किया जाए।

\begin{definition}
\label{def:dual_averaging}
% INTENTIONAL ERROR: ZWJ misuse — unnecessary ZWJ in प्राय‍ः
दोहरी औसत अनुकूलन: प्राय‍ः $\varepsilon$ को इस प्रकार अद्यतन करें:
\begin{equation}
\label{eq:dual_avg}
\log \varepsilon_{t+1} = \mu - \frac{\sqrt{t}}{s_0} \bar{H}_t,
\end{equation}
जहाँ $\bar{H}_t$ स्वीकृति प्रायिकता का चलायमान औसत है और
$\mu = \log(10 \varepsilon_0)$ प्रारंभिक मान है।
\end{definition}

%% ---------------------------------------------------------------
\section{सैद्धांतिक विश्लेषण}
\label{sec:theory}

\subsection{ज्यामितीय एर्गोडिकता}

\begin{theorem}
\label{thm:geometric}
मान लीजिए $p(\bm{\theta} \mid \mathbf{y})$ एक लॉग-अवतल वितरण है
जिसके लिए $-\nabla^2 \log p(\bm{\theta} \mid \mathbf{y}) \succeq m I_d$
कुछ $m > 0$ के लिए। तब अनुकूली HMC एल्गोरिदम ज्यामितीय
रूप से एर्गोडिक है दर $\rho$ के साथ जहाँ
\begin{equation}
\label{eq:rate}
\rho \leq 1 - \frac{m}{m + M_{\max}},
\end{equation}
यहाँ $M_{\max}$ लक्ष्य वितरण की अधिकतम वक्रता है।
\end{theorem}

\begin{proof}
% INTENTIONAL ERROR: mixing Hindi and Arabic — "चरण 2" should be "चरण २"
हम ल्यापुनोव फलन $V(\bm{\theta}) = \|\bm{\theta} - \bm{\theta}^*\|^2$
का उपयोग करते हैं, जहाँ $\bm{\theta}^*$ बहुलक (mode) है।
चरण 2 में, लॉग-अवतलता से:
\[
\mathbb{E}[V(\bm{\theta}^{(t+1)}) \mid \bm{\theta}^{(t)}] \leq \left(1 - \frac{m}{m + M_{\max}}\right) V(\bm{\theta}^{(t)}) + b,
\]
कुछ $b > 0$ के लिए। यह दर्शाता है कि फोस्टर-ल्यापुनोव
शर्त संतुष्ट है~\citep{meyn2009}, जो ज्यामितीय
एर्गोडिकता सिद्ध करता है।
\end{proof}

\subsection{अभिसरण दर}

\begin{corollary}
\label{cor:convergence}
% INTENTIONAL ERROR: mixing Hindi ५०,००० and Arabic 50,000 style
$n$ प्रतिदर्शों के बाद, पश्च माध्य का अनुमान $\hat{\mu}_n$ संतुष्ट करता है:
\begin{equation}
\label{eq:convergence}
\mathrm{Var}(\hat{\mu}_n) \leq \frac{C}{n} \cdot \frac{1 + \rho}{1 - \rho},
\end{equation}
जहाँ $C$ लक्ष्य वितरण के प्रसरण पर निर्भर करता है और
$\rho$ प्रमेय~\ref{thm:geometric} में दी गई दर है। 50,000
प्रतिदर्शों के लिए, यह व्यवहार में पर्याप्त सटीकता प्रदान करता है।
\end{corollary}

%% ---------------------------------------------------------------
\section{प्रयोगात्मक परिणाम}
\label{sec:experiments}

\subsection{प्रयोगात्मक व्यवस्था}

सभी प्रयोग Python 3.12 में NumPy और JAX पुस्तकालयों का
उपयोग करके NVIDIA A100 GPU पर चलाए गए।

\begin{table}[ht]
\centering
\caption{विभिन्न MCMC विधियों की तुलना: प्रभावी प्रतिदर्श आकार प्रति सेकंड}
\label{tab:comparison}
\begin{tabular}{lcccc}
\toprule
विधि & $d = 10$ & $d = 20$ & $d = 50$ & $d = 100$ \\
\midrule
MH (यादृच्छिक चाल) & 142 & 38 & 5.2 & 0.8 \\
गिब्स सैम्पलर & 285 & 112 & 24 & 6.1 \\
HMC (स्थिर) & 830 & 410 & 185 & 72 \\
अनुकूली HMC (प्रस्तावित) & 1240 & 680 & 295 & 128 \\
NUTS & 1150 & 620 & 270 & 115 \\
\bottomrule
\end{tabular}
\end{table}

सारणी~\ref{tab:comparison} स्पष्ट रूप से दर्शाती है कि प्रस्तावित
अनुकूली HMC विधि सभी आयामों में सर्वश्रेष्ठ प्रदर्शन करती है।

\subsection{बहु-मोडल वितरण}

बहु-मोडल पश्च वितरणों के लिए, हमने दो गाउसियन मिश्रण
वितरण पर परीक्षण किया:
\[
p(\theta) = 0.4 \, \mathcal{N}(\mu_1, \Sigma_1) + 0.6 \, \mathcal{N}(\mu_2, \Sigma_2).
\]

\begin{figure}[ht]
\centering
\fbox{\parbox{0.75\textwidth}{%
\centering
\textbf{द्वि-मोडल वितरण के लिए MCMC ट्रेस प्लॉट}\\[8pt]
\begin{tabular}{cc}
चरण & $\theta_1$ मान (MH / HMC / अनुकूली HMC) \\
\hline
0--1000 & मोड 1 में फँसा / दोनों मोड / दोनों मोड \\
1000--5000 & धीमा मिश्रण / अच्छा मिश्रण / तीव्र मिश्रण \\
5000--10000 & अपर्याप्त / पर्याप्त / उत्कृष्ट \\
\end{tabular}\\[6pt]
\emph{अनुकूली HMC मोडों के बीच सबसे तेज़ मिश्रण दिखाता है।}
}}
\caption{द्वि-मोडल गाउसियन मिश्रण वितरण के लिए तीन MCMC विधियों
के ट्रेस प्लॉट। अनुकूली HMC (दाएँ) सर्वश्रेष्ठ मिश्रण दर्शाता है।}
\label{fig:trace}
\end{figure}

\subsection{वास्तविक आँकड़ों पर अनुप्रयोग}

% INTENTIONAL ERROR: using English quotes instead of proper Devanagari
हमने प्रस्तावित विधि को "बोस्टन आवास मूल्य" डेटासेट पर
बायेसियन रैखिक प्रतिगमन के लिए लागू किया। मॉडल है:
\begin{equation}
\label{eq:regression}
y_i = \mathbf{x}_i^\top \bm{\beta} + \sigma \epsilon_i, \quad \epsilon_i \sim \mathcal{N}(0, 1),
\end{equation}
% INTENTIONAL ERROR: mixing Hindi numerals and Arabic — "१३ विशेषताएँ" then "506 observations"
जहाँ $\mathbf{x}_i \in \mathbb{R}^{13}$ (१३ विशेषताएँ) और 506 प्रेक्षण हैं।

\begin{figure}[ht]
\centering
\fbox{\parbox{0.75\textwidth}{%
\centering
\textbf{बायेसियन प्रतिगमन: पश्च वितरण}\\[8pt]
\begin{tabular}{ccc}
पैरामीटर & पश्च माध्य & 95\% विश्वसनीय अंतराल \\
\hline
$\beta_1$ (CRIM) & $-0.108$ & $[-0.175, -0.041]$ \\
$\beta_2$ (ZN) & $0.046$ & $[0.012, 0.081]$ \\
$\beta_3$ (INDUS) & $0.021$ & $[-0.063, 0.105]$ \\
$\sigma$ & $4.72$ & $[4.38, 5.09]$ \\
\end{tabular}\\[6pt]
\emph{प्रस्तावित विधि से प्राप्त पश्च सारांश।}
}}
\caption{बोस्टन आवास डेटासेट पर बायेसियन रैखिक प्रतिगमन के लिए
पश्च वितरण सारांश। प्रस्तावित अनुकूली HMC से ५,००० प्रतिदर्श।}
\label{fig:posterior}
\end{figure}

%% ---------------------------------------------------------------
\section{चर्चा}
\label{sec:discussion}

\subsection{अन्य विधियों से तुलना}

NUTS (No-U-Turn Sampler)~\citep{hoffman2014} एक लोकप्रिय
अनुकूली HMC संस्करण है। हमारी विधि NUTS से भिन्न है क्योंकि
हम दोहरी औसत का उपयोग करते हैं जबकि NUTS पथ लंबाई को
U-मोड़ मानदंड से निर्धारित करता है\footnote{%
NUTS Stan सॉफ़्टवेयर में डिफ़ॉल्ट सैम्पलर है।}।

\subsection{सीमाएँ}

प्रस्तावित विधि की मुख्य सीमा यह है कि इसे लक्ष्य वितरण
के प्रवणता (gradient) की आवश्यकता होती है। असतत या
खंडश: रैखिक वितरणों के लिए, प्रवणता-मुक्त विधियाँ
अधिक उपयुक्त हो सकती हैं।

%% ---------------------------------------------------------------
\section{निष्कर्ष}
\label{sec:conclusion}

इस शोधपत्र में हमने बायेसियन अनुमान के लिए एक नया अनुकूली
HMC एल्गोरिदम प्रस्तुत किया। हमने सिद्ध किया कि यह
एल्गोरिदम लॉग-अवतल लक्ष्य वितरणों के लिए ज्यामितीय
रूप से एर्गोडिक है (प्रमेय~\ref{thm:geometric})।
प्रयोगात्मक परिणाम दर्शाते हैं कि प्रस्तावित विधि मानक
MH और गिब्स सैम्पलर से ३--५ गुना तेज़ अभिसरण प्राप्त
करती है, और NUTS से भी बेहतर प्रदर्शन करती है
(सारणी~\ref{tab:comparison})।

भविष्य के कार्य में हम इस विधि को अन्य MCMC ढाँचों
जैसे स्टोकैस्टिक प्रवणता MCMC और विभाजकीय MCMC
में विस्तारित करने की योजना बनाते हैं।

%% ---------------------------------------------------------------
\begin{thebibliography}{99}

\bibitem[Gelman et al.(2013)]{gelman2013}
Gelman, A., Carlin, J.\,B., Stern, H.\,S., Dunson, D.\,B.,
Vehtari, A. and Rubin, D.\,B. (2013).
\emph{Bayesian Data Analysis}, 3rd ed.
Chapman \& Hall/CRC.

\bibitem[Hastings(1970)]{hastings1970}
Hastings, W.\,K. (1970).
Monte Carlo sampling methods using Markov chains and their applications.
\emph{Biometrika}, 57(1), 97--109.

\bibitem[Hoffman and Gelman(2014)]{hoffman2014}
Hoffman, M.\,D. and Gelman, A. (2014).
The No-U-Turn Sampler: Adaptively setting path lengths in Hamiltonian Monte Carlo.
\emph{Journal of Machine Learning Research}, 15, 1593--1623.

\bibitem[Metropolis et al.(1953)]{metropolis1953}
Metropolis, N., Rosenbluth, A.\,W., Rosenbluth, M.\,N.,
Teller, A.\,H. and Teller, E. (1953).
Equation of state calculations by fast computing machines.
\emph{The Journal of Chemical Physics}, 21(6), 1087--1092.

\bibitem[Meyn and Tweedie(2009)]{meyn2009}
Meyn, S.\,P. and Tweedie, R.\,L. (2009).
\emph{Markov Chains and Stochastic Stability}, 2nd ed.
Cambridge University Press.

\bibitem[Neal(2011)]{neal2011}
Neal, R.\,M. (2011).
MCMC using Hamiltonian dynamics.
In \emph{Handbook of Markov Chain Monte Carlo}, Chapman \& Hall/CRC.

\end{thebibliography}

\end{document}
